% Generated from locale/tr/content/incompleteness/representability-in-q/introduction.tex; do not hand-edit.


\olfileid[tr]{inc}{req}{int}
\olsection{Giriş}

Eksiklik teoremleri, hesaplanabilir fonksiyonlar hakkındaki temel olguların
dile getirilip ispatlanabildiği kuramlara uygulanır. Böyle, son derece küçük
bir kuramı betimleyeceğiz; buna ``$\Th{Q}$'' (ya da Raphael Robinson'a
atıfla bazen ``Robinson $Q$'su'') denir. Bir fonksiyonun $\Th{Q}$ içinde
\emph{temsil edilebilir} olmasının ne anlama geldiğini söyleyecek ve ardından
şunu ispatlayacağız:
\begin{quote}
  Bir fonksiyon $\Th{Q}$ içinde ancak ve ancak hesaplanabilirse temsil
  edilebilir.
\end{quote}
Bu, her şeyden önce bize başka bir hesaplanabilirlik modeli sağlar. Fakat
durma problemini ona indirgeyerek $\Setabs{!A}{\Th{Q} \Proves !A}$ kümesinin
karar verilebilir olmadığını göstermek için de bu sonucu kullanacağız.
Bitirdiğimizde bundan çok daha güçlü şeyler ispatlamış olacağız.

$\Th{Q}$'nun dili aritmetik dilidir; $\Th{Q}$ aşağıdaki aksiyomlardan oluşur
(bunlar !!{identity} içeren birinci derece mantığın öteki aksiyom ve
kurallarıyla birlikte kullanılacaktır):
\begin{align*}
& \lforall[x][\lforall[y][(\eq[x'][y'] \lif \eq[x][y])]] \tag{$!Q_1$}\\
& \lforall[x][\eq/[\Obj 0][x']] \tag{$!Q_2$}\\
& \lforall[x][(\eq[x][\Obj 0] \lor \lexists[y][\eq[x][y']])] \tag{$!Q_3$}\\
& \lforall[x][\eq[(x + \Obj 0)][x]] \tag{$!Q_4$}\\
& \lforall[x][\lforall[y][\eq[(x + y')][(x + y)']]] \tag{$!Q_5$}\\
& \lforall[x][\eq[(x \times \Obj 0)][\Obj 0]] \tag{$!Q_6$}\\
& \lforall[x][\lforall[y][\eq[(x \times y')][((x \times y) + x)]]] \tag{$!Q_7$}\\
& \lforall[x][\lforall[y][(x < y \liff \lexists[z][\eq[(z' + x)][y]])]] \tag{$!Q_8$}
\end{align*}
Her $n$ doğal sayısı için $\num{n}$ sayısalını, toplam $n$ tane kesme imi
bulunan $\Obj{0}^{\prime\prime\ldots\prime}$ terimi olarak tanımlayın.
Dolayısıyla $\num{0}$ tek başına !!{constant}~$\Obj{0}$, $\num{1}$
$\Obj{0}'$, $\num{2}$ ise $\Obj{0}''$'dır; böyle devam eder.

Bir aritmetik kuramı olarak $\Th{Q}$ \emph{son derece} zayıftır; örneğin
$\lforall[x][\eq/[x][x']]$ veya $\lforall[x][\lforall[y][(x + y) = (y +
    x)]]$ gibi çok yalın olguları bile ispatlayamazsınız. Fakat
$\Th{Q}$'yu böylesine ilginç kılan nedenlerin çoğunun onun bu kadar zayıf
\emph{olması} olduğunu göreceğiz. Gerçekte eksiklik teoreminin geçerli
olmasına ancak yetecek güçtedir. $\Th{Q}$'nun ilginç olmasının başka bir
nedeni, \emph{sonlu} bir aksiyomlar kümesine sahip olmasıdır.

$\Th{Q}$'dan daha güçlü olan ve \emph{Peano aritmetiği} $\Th{PA}$ denen
kuram, $\Th{Q}$'ya tümevarım şeması eklenerek elde edilir:
\[
(!A(\Obj 0) \land \lforall[x][(!A(x) \lif !A(x'))]) \lif \lforall[x][!A(x)]
\]
burada $!A(x)$ herhangi bir formüldür. $!A(x)$, $x$ dışında serbest
!!{variable}s içeriyorsa bunların hepsini bağlamak için başa evrensel
niceleyiciler ekleriz (böylece tümevarım şemasının karşılık gelen örneği
!!a{sentence} olur). Örneğin $!A(x,y)$ serbest olarak !!{variable}~$y$'yi
de içeriyorsa karşılık gelen örnek şudur:
\[
\lforall[y][((!A(\Obj 0) \land \lforall[x][(!A(x) \lif !A(x'))]) \lif
  \lforall[x][!A(x)])]
\]
Tümevarım şemasının örneklerini kullanarak $\Th{PA}$'nın aksiyomlarından,
$\Th{Q}$'nun aksiyomlarından çıkarılabileceğinden çok daha fazlası
ispatlanabilir. Gerçekte doğal sayılar hakkında Peano aritmetiğinde
ispatlanamayan ``doğal'' ifadeler bulmak epey emek gerektirir!{}

\begin{defn}
\ollabel{defn:representable-fn}
  Doğal sayılardan doğal sayılara giden $f(x_0,\ldots,x_k)$ fonksiyonuna,
  $f(n_0,\dots,n_k)=m$ olduğunda $\Th{Q}$ aşağıdakileri ispatlayacak
  biçimde bir $!A_f(x_0,\dots,x_k,y)$ formülü varsa
  {\em $\Th{Q}$ içinde temsil edilebilir} denir:
\begin{enumerate}
\item\ollabel{defn:rep:a} $!A_f(\num{n_0}, \dots, \num{n_k}, \num{m})$
\item\ollabel{defn:rep:b} $\lforall[y][(!A_f(\num{n_0}, \dots,
\num{n_k}, y) \lif \num{m} = y)]$.
\end{enumerate}
\end{defn}

Tanımı başka biçimlerde de ifade edebiliriz; örneğin $\Th{Q}$'nun
$\lforall[y][(!A_f(\num{n_0}, \dots, \num{n_k}, y) \liff
\eq[y][\num{m}])]$ tümcesini ispatlamasını denk bir koşul olarak
isteyebilirdik.

\begin{thm}
\ollabel{thm:representable-iff-comp}
Bir fonksiyon $\Th{Q}$ içinde ancak ve ancak hesaplanabilirse temsil
edilebilir.
\end{thm}

Teoremin ispatlanacak iki yönü vardır. Soldan sağa yön, sözdizimin
aritmetikleştirilmesi hazır olduğunda oldukça doğrudandır. Öteki yön daha
fazla emek gerektirir. Temel düşünce şudur: ``hesaplanabilir''i kesinleştirmenin
bir yolu olarak ``genel özyinelemeli''yi seçer ve her genel özyinelemeli
fonksiyonun $\Th{Q}$ içinde temsil edilebilir olduğunu gösteririz. Bir
fonksiyonun; $\Zero$, ardıl fonksiyonu~$\Succ$ ve izdüşüm
fonksiyonları~$\Proj{n}{i}$'den bileşke, ilkel özyineleme ve düzenli
minimizasyon kullanılarak tanımlanabiliyorsa genel özyinelemeli olduğunu
anımsayın. Bu nedenle her genel özyinelemeli fonksiyonun $\Th{Q}$ içinde
temsil edilebilir olduğunu göstermenin bir yolu, temel fonksiyonların temsil
edilebilir olduğunu ve bazı fonksiyonlar temsil edilebilirse bunlardan
bileşke, ilkel özyineleme ve düzenli minimizasyon kullanılarak tanımlanan
fonksiyonların da temsil edilebildiğini göstermektir. Başka bir deyişle,
temel fonksiyonların temsil edilebilir olduğunu ve temsil edilebilir
fonksiyonların bileşke, ilkel özyineleme ve düzenli minimizasyon ``altında
kapalı'' olduğunu gösterebiliriz. Bu, her genel özyinelemeli fonksiyonun
temsil edilebilir olmasını güvenceye alır.

Temsil edilebilir fonksiyonların ilkel özyineleme altında kapalı olduğunu
göstereceğimiz adımın zor olduğu ortaya çıkar. Bu adımdan kaçınmak için önce
ilkel özyineleme olmadan da çalışabileceğimizi gösteririz. Yani her genel
özyinelemeli fonksiyonun yalnızca temel fonksiyonlardan bileşke ve düzenli
minimizasyon kullanılarak tanımlanabileceğini gösteririz. Bunu yapmak için
ilkel özyinelemenin gerçekte belirli bir düzenli minimizasyonla yapılabildiğini
gösteririz. Ancak bunun işe yaraması için birtakım ek temel fonksiyonlar
katmalıyız: toplama, çarpma ve özdeşlik bağıntısının karakteristik
fonksiyonu~$\Char{=}$. Ardından \emph{bu} temel fonksiyonların hepsinin
$\Th{Q}$ içinde temsil edilebilir olduğunu ve temsil edilebilir fonksiyonların
bileşke ile düzenli minimizasyon altında kapalı olduğunu göstererek teoremi
ispatlayabiliriz.

